\chapter{So sánh bốn phương pháp và đánh đổi triển khai}
\label{chap:comparison}
\section{Điều kiện đối chiếu và nguồn số liệu}
So sánh dùng corpus 199.965 bài, taxonomy 20 lớp và cùng phép chia phân tầng 70/15/15, seed 42. Đối chiếu trực tiếp xác nhận mảng chỉ số train, validation và test của bốn notebook giống nhau. Trên test, toàn bộ 29.995 URL, nhãn tham chiếu và thứ tự mẫu khớp giữa ba đối chứng và metadata PhoBERT. Số lượng bằng nhau hoặc chuỗi fingerprint đơn lẻ chưa đủ cho kết luận này: PhoBERT và các đối chứng dùng cách tạo signature khác nhau.

Các số được đọc từ metrics, scores/logits, predictions, ma trận và classification report đã xuất. Tái tính từ từng hàng dự đoán xác nhận Accuracy, Macro-F1, Weighted-F1 và tổng support. Mỗi model đi vào test với duy nhất cấu hình/checkpoint đã chọn trên validation. SVM/LR được tìm trong 12 cấu hình mỗi model, BiLSTM trong ba cấu hình, PhoBERT dùng cấu hình fine-tuning ở chương~\ref{chap:training}. Ngân sách lựa chọn không bằng nhau; đây là so sánh các pipeline thực thi, không cô lập ưu thế kiến trúc.

\section{Chất lượng phân loại trên cùng tập test}
\DataTable{tables/comparison_primary.tex}{Chất lượng bốn model với quyết định argmax score/logit; chưa dùng hệ số điều chỉnh lớp của PhoBERT. Các chỉ số phần trăm tính trên cùng 29.995 bài test.}
PhoBERT đạt Macro-F1 87,22\%, SVM 85,33\%, LR 84,98\% và BiLSTM 84,60\%. Khoảng cách PhoBERT--SVM là 1,89 điểm phần trăm Macro-F1 và 1,17 điểm Accuracy, tương ứng 350 bài đúng ròng. SVM có 26.970 bài đúng, so với 26.714 của LR và 26.595 của BiLSTM. Những giá trị này là kết quả corpus hiện tại, không lấy từ bảng nghiên cứu ở chương~\ref{chap:related}.

Temperature dương không đổi argmax ở cả bốn model. Riêng hệ số lớp của PhoBERT thay nhãn sau temperature, nâng Accuracy từ 91,08\% lên 91,13\% và Macro-F1 từ 87,22\% lên 87,25\%, đạt 27.333 bài đúng. Trạng thái này phục vụ ứng dụng ở chương~\ref{chap:application}; nó được phân tích riêng trong chương~\ref{chap:results}, không dùng thay dòng argmax trong đối chứng chính.

Kết quả cho thấy TF--IDF toàn văn đã cung cấp đối chứng mạnh. BiLSTM học từ đầu chưa vượt hai mô hình tuyến tính trong cấu hình đã khảo sát. Điều đó không chứng minh thứ tự từ vô ích, LSTM nói chung kém hơn SVM, hay toàn bộ mức tăng của PhoBERT đến từ pretraining. Cách cắt, stopword, lặp tiêu đề, trọng số lớp và ngân sách tối ưu đều là biến đồng thời thay đổi.

\section{Chênh lệch theo chuyên mục và lớp ít mẫu}
\DataTable{tables/comparison_per_class.tex}{F1 (\%) theo 20 chuyên mục với argmax, cùng support test; DTTS là Dân tộc thiểu số. Giá trị làm tròn có thể che chênh lệch rất nhỏ.}
Thể thao, Thế giới và Ô tô -- Xe máy đạt F1 cao ở cả bốn pipeline. SVM và PhoBERT đều khoảng 98,78\% ở Thể thao; khác biệt tại chữ số nhỏ không đủ để coi một model tốt hơn có ý nghĩa theo lớp này. Ngược lại, Dân tộc thiểu số và miền núi đạt 69,38\% với PhoBERT, cao hơn mức 58,06--59,53\% của ba đối chứng. Nội dung chuyên đề đạt 79,74\% với PhoBERT, so với 74,18--76,57\%. Lợi thế tổng thể vì vậy không phân bố đều trên mọi chuyên mục.

\subsection{Nông thôn mới: đánh đổi precision và recall}
\DataTable{tables/comparison_rare.tex}{Nông thôn mới: precision, recall và F1 của bốn model trước điều chỉnh nhãn; support là 146 bài.}
BiLSTM đạt F1 48,24\%, LR 47,06\%, PhoBERT argmax 45,57\% và SVM 39,53\%. Ngay cả PhoBERT sau điều chỉnh lớp đạt 46,39\%, vẫn thấp hơn hai giá trị đầu trong mẫu test này. SVM nhận đúng 51/146 bài, LR 80, BiLSTM 89 và PhoBERT argmax 72. Recall cao hơn của BiLSTM đi kèm precision chỉ 39,91\%; đây không phải cải thiện đồng thời mọi tiêu chí.

Ma trận cho thấy SVM chuyển 86 bài Nông thôn mới sang Nội dung chuyên đề. Ở chiều ngược, LR chuyển 106 bài Nội dung chuyên đề sang Nông thôn mới, BiLSTM chuyển 128 bài. Những lỗi từ lớp lớn góp phần làm precision của lớp hiếm thấp. Trọng số loss, số mẫu và đặc điểm nội dung cùng ảnh hưởng ranh giới; chưa có ablation để xác định nguyên nhân riêng. Với 146 mẫu, thứ hạng F1 theo lớp cần được kiểm chứng trên nhiều seed và một tập gán nhãn được rà soát.

\subsection{Ma trận chuẩn hóa và phân bố F1}
Các hình ma trận dùng cùng ID của bảng taxonomy. Mỗi hàng được chuẩn hóa theo support lớp, nên đọc như phân bố dự đoán khi biết nhãn tham chiếu. Hình SVM/LR/BiLSTM dùng đơn vị phần trăm, còn ma trận chuẩn hóa PhoBERT ở chương~\ref{chap:results} dùng tỷ lệ 0--1; hai thang tương đương sau nhân 100, không được coi màu là một thang số chung chưa kiểm tra.
\PortraitFigure{assets/baselines/svm/evaluation/confusion_matrix_normalized.png}{Ma trận SVM chuẩn hóa theo hàng, đơn vị \%; một hình riêng trên A4.}
\PortraitFigure{assets/baselines/logistic/evaluation/confusion_matrix_normalized.png}{Ma trận Logistic Regression chuẩn hóa theo hàng, đơn vị \%; cùng taxonomy và tập test.}
\PortraitFigure{assets/baselines/lstm/evaluation/confusion_matrix_normalized.png}{Ma trận BiLSTM chuẩn hóa theo hàng, đơn vị \%; không ghép với ma trận model khác.}
\PortraitFigure{assets/baselines/svm/evaluation/per_class_f1.png}{F1 theo chuyên mục của SVM, trên 29.995 bài test.}
\PortraitFigure{assets/baselines/logistic/evaluation/per_class_f1.png}{F1 theo chuyên mục của Logistic Regression, trên cùng tập test.}
\PortraitFigure{assets/baselines/lstm/evaluation/per_class_f1.png}{F1 theo chuyên mục của BiLSTM; Nông thôn mới không đồng thứ hạng với điểm tổng.}

\section{Phân tích ghép cặp và bất định của so sánh}
\label{sec:paired_comparison}
Hai model dự đoán cùng bài nên cần xét chênh lệch ghép cặp, không chỉ nhìn khoảng tin cậy riêng. Phân tích tại báo cáo lấy 1.000 bootstrap IID, seed 42 bằng \texttt{numpy.random.default\_rng}, mỗi lần lấy có hoàn lại 29.995 vị trí. Cùng một mảng vị trí được áp dụng cho cả bốn model; F1 luôn tính trên đủ 20 nhãn, với mẫu số bằng không nhận 0. Với cặp $A,B$:
\begin{equation}
\Delta^{(b)}=\operatorname{MacroF1}(y_{I_b},\hat y^A_{I_b})-
\operatorname{MacroF1}(y_{I_b},\hat y^B_{I_b}).
\end{equation}
Khoảng percentile 95\% lấy phân vị 2,5 và 97,5 của $\Delta^{(b)}$. Đây là hậu phân tích từ dự đoán đã xuất, không chạy lại notebook hoặc huấn luyện. Các model và tham số giữ cố định; kết quả khác khoảng riêng 500 lần của từng notebook vì mục tiêu là chênh lệch có bảo toàn ghép cặp.
\DataTable{tables/comparison_paired.tex}{Chênh lệch Macro-F1 ghép cặp với argmax; khoảng tin cậy là từng cặp cho model cố định, không phải khoảng đồng thời của toàn bộ so sánh.}
Các khoảng PhoBERT--SVM/LR/BiLSTM đều nằm trên 0 trong thiết lập này. Khoảng SVM--LR là $[-0,09;0,75]$ điểm phần trăm, LR--BiLSTM là $[-0,17;0,95]$, đều chứa 0. Vì vậy, thứ hạng điểm mẫu giữa các đối chứng gần nhau chưa đủ cho kết luận mạnh về thứ hạng tổng quát. Bootstrap không đo biến thiên do seed huấn luyện, không xử lý phụ thuộc giữa bài cùng sự kiện và không thay thế test theo thời gian hoặc nguồn.

PhoBERT đúng/SVM sai ở 1.439 bài, còn SVM đúng/PhoBERT sai ở 1.089 bài. Chênh lệch ròng 350 bài không có nghĩa lỗi SVM là một tập con của lỗi PhoBERT. Hai pipeline có vùng mạnh khác nhau; việc tận dụng sự khác biệt bằng ensemble là một giả thuyết, chưa phải kết quả đo trong báo cáo.

\section{Lỗi chung và ranh giới nhãn biên tập}
\DataTable{tables/comparison_shared_errors.tex}{Số model dự đoán đúng từng bài trong bốn quyết định argmax; tổng bằng 29.995.}
Cả bốn model cùng đúng ở 24.677 bài và cùng sai ở 1.251 bài. Trong nhóm cùng sai, Nội dung chuyên đề có 186 bài, Kinh doanh 169 và Đời sống 156; Nông thôn mới có 24/146 bài, Dân tộc -- Tôn giáo 83/555 bài. Số tuyệt đối và tỷ lệ trong lớp thể hiện hai mặt khác nhau của độ khó, không được dùng thay nhau.

Ba ví dụ chung có thể truy vết từ artifact: ``Google tích hợp AI Gemini vào hệ sinh thái sản phẩm'' có nhãn Giáo dục nhưng cả bốn dự đoán Công nghệ; ``The Sonata -- sống tận hưởng tại `tọa độ quốc tế' bên sông Hàn'' có nhãn Kinh doanh nhưng cả bốn chọn Bất động sản; ``Ô tô mới có thể được miễn đăng kiểm lần đầu từ năm 2023'' có nhãn Thời sự nhưng cả bốn chọn Ô tô -- Xe máy. URL tương ứng là:
\begin{itemize}
\item \href{https://vietnamnet.vn/google-tich-hop-ai-gemini-vao-he-sinh-thai-san-pham-2281849.html}{VietNamNet, bài 2281849: Google tích hợp AI Gemini};
\item \href{https://vietnamnet.vn/the-sonata-song-tan-huong-tai-toa-do-quoc-te-ben-song-han-2348982.html}{VietNamNet, bài 2348982: The Sonata bên sông Hàn};
\item \href{https://vietnamnet.vn/o-to-moi-co-the-duoc-mien-dang-kiem-lan-dau-tu-nam-2023-2094970.html}{VietNamNet, bài 2094970: Miễn đăng kiểm lần đầu}.
\end{itemize}
Những trường hợp này cho thấy nhãn chuyên mục có thể phụ thuộc cách tổ chức biên tập hơn tên chủ đề dễ nhận thấy. Không tự đổi nhãn tham chiếu vì nhiều model đồng ý với nhau. Cần rà soát bài gốc, quy tắc taxonomy và bằng chứng chuyên mục trước khi kết luận nhãn sai hoặc model sai theo nghĩa ngữ nghĩa.

\section{Calibration: chất lượng nhãn và độ tin cậy}
\DataTable{tables/comparison_calibration.tex}{NLL và ECE 15 bins trước/sau temperature; confidence là xác suất cực đại trước hệ số quyết định lớp.}
SVM giảm NLL từ 1,2241 xuống 0,3540 và ECE từ 0,5457 xuống 0,0207 với $T=0,2718$. Đây là hiệu chỉnh thang margin: softmax trực tiếp từ margin chưa phải xác suất đã hiệu chỉnh. LR dùng $T=0,7950$; BiLSTM và PhoBERT dùng $T>1$, làm phân bố bớt tập trung. Không suy ra chỉ từ dấu $T-1$ rằng mọi lớp hoặc mọi mức confidence đều có cùng loại sai lệch.

Sau temperature, BiLSTM/LR có ECE khoảng 0,0115/0,0119, thấp hơn PhoBERT 0,0251; PhoBERT có NLL thấp nhất, 0,3220. ECE là tổng sai lệch confidence--accuracy theo các bin, còn NLL xét xác suất nhãn tham chiếu của từng mẫu. Model có ECE thấp nhất không đồng nghĩa model có F1 cao nhất hoặc xác suất tối ưu theo mọi tiêu chí. Không dùng ngưỡng confidence của app PhoBERT cho một model khác mà chưa kiểm chứng lại định tuyến.
\ReportFigure{assets/baselines/svm/evaluation/reliability_before.png}{Reliability của softmax margin SVM trước temperature; thang score chưa được hiệu chỉnh.}
\ReportFigure{assets/baselines/svm/evaluation/reliability_temperature.png}{Reliability của SVM sau temperature được fit trên validation.}
\ReportFigure{assets/baselines/logistic/evaluation/reliability_temperature.png}{Reliability của Logistic Regression sau temperature, dùng 15 bin confidence.}
\ReportFigure{assets/baselines/lstm/evaluation/reliability_temperature.png}{Reliability của BiLSTM sau temperature; ECE thấp không thay thế chỉ số nhãn.}

\section{Chi phí tìm cấu hình, huấn luyện và suy luận}
\DataTable{tables/comparison_timing.tex}{Thời gian được ghi trong artifact của từng pipeline; dấu -- ở PhoBERT nghĩa không có lượt tìm nhiều cấu hình tương ứng. Phạm vi timer cần đọc cùng diễn giải.}
SVM fit classifier được chọn trong 27,319 giây, LR trong 239,531 giây, đều dùng CPU với bốn luồng. Tổng tìm 12 cấu hình là 758,245 và 2.370,248 giây, gồm tạo TF--IDF và các lượt tối ưu. BiLSTM train cấu hình chọn trong 1.779,429 giây (29,66 phút), tìm ba cấu hình trong 4.473,369 giây (74,56 phút). PhoBERT fine-tuning bảy epoch mất 12.429,912 giây (207,17 phút), không tính chi phí tiền huấn luyện mô hình nền. Timer fit SVM/LR chỉ đo tối ưu classifier; timer huấn luyện neural gồm vòng train, validation theo epoch và lưu checkpoint, nên không phải các hạng chi phí hoàn toàn đồng nhất.

Phạm vi suy luận khác nhau. SVM/LR gồm TF--IDF transform và decision function trên CPU, không gồm phân đoạn tiếng Việt. BiLSTM gồm DataLoader, chuyển tensor sang GPU và forward, không gồm phân đoạn/mã hóa; batch test là 128. Timer test PhoBERT là 106,322 giây trên GPU, không gồm nạp model; batch test là 64 và input đã mã hóa. Các thời gian 12,689/13,125/9,233 giây của ba đối chứng phản ánh workload được ghi, không phải latency từ URL đến nhãn trong app. Chưa có benchmark serving CPU đồng nhất để xếp hạng độ trễ triển khai.

BiLSTM có 10,34 triệu tham số, PhoBERT 135,01 triệu; hai mô hình tuyến tính lưu hệ số, từ vựng và IDF, không có encoder lớn. Artifact BiLSTM ghi peak allocated 0,564 GiB trong huấn luyện cấu hình được chọn, metadata PhoBERT ghi 5,388 GiB ở lượt đánh giá. Hai mốc đo khác giai đoạn; chưa có peak RAM/VRAM toàn quy trình đồng nhất cho cả bốn model. Không dùng dung lượng file nén hoặc bộ nhớ CUDA reserved thay cho bộ nhớ thực thi được đo.

\section{Trả lời câu hỏi nghiên cứu và giới hạn kết luận}
\textbf{RQ1:} trên snapshot và phép chia hiện tại, PhoBERT argmax có Macro-F1 cao nhất; SVM là đối chứng tuyến tính mạnh với chi phí fit thấp. Bất định ghép cặp ủng hộ khoảng cách PhoBERT với ba đối chứng khi giữ model cố định, nhưng không xác nhận thứ hạng SVM--LR--BiLSTM ở mọi lần train.

\textbf{RQ2:} mức tăng khoảng 1,89 điểm Macro-F1 so với SVM đi kèm fine-tuning GPU tốn thời gian hơn và encoder lớn hơn. Các timer cho phép nhận diện đánh đổi trong thực nghiệm, chưa tạo một phép đo latency/bộ nhớ serving hoàn toàn đồng nhất. Lựa chọn triển khai còn phụ thuộc tài nguyên, lưu lượng và yêu cầu rà soát, không chỉ một F1 tổng thể.

\textbf{RQ3:} khác biệt theo lớp có thật trong mẫu: BiLSTM/LR tốt hơn ở Nông thôn mới, PhoBERT tốt hơn rõ hơn ở Dân tộc thiểu số và miền núi, còn các lớp chủ đề tập trung dễ đạt điểm cao với TF--IDF. Lỗi chung ở ranh giới nhãn biên tập cho thấy cần đầu tư kiểm toán taxonomy và mẫu khó, không chỉ tăng độ phức tạp model.

Tất cả kết luận trên thuộc một nguồn báo, một seed và các pipeline đã công bố. Không có test tương lai/ngoài nguồn, thực nghiệm nhiều seed hoặc ablation cô lập pretraining, lặp tiêu đề, stopword và chiều dài. Kết quả không xác lập PhoBERT luôn tốt nhất trên tin tức tiếng Việt; ứng dụng minh họa tiếp tục dùng gói PhoBERT và chính sách hiệu chỉnh đã được kiểm chứng.
